\section{Administración en tiempo de ejecución del historial de mensajes}

\subsection{Administración del Context en interacciones de múltiples turnos}

Después de los 4 mensajes preinyectados, Codex entra en la etapa real de interacción de múltiples turnos entre el usuario y el Agent:

\begin{enumerate}
    \item El usuario emite una Query real
    \item El Agent Loop comienza a ejecutar una interacción de múltiples turnos (multi-turn)
    \item Durante la interacción se incluyen llamadas a herramientas (Function Call) y sus resultados de retorno
    \item Termina un Turn
    \item El usuario emite una nueva Query, con lo que comienza el siguiente Turn
\end{enumerate}

\begin{importantbox}{La función de la instrucción Compact}
Cuando el historial de mensajes se acumula en exceso, se puede ejecutar la instrucción \texttt{compact} en Codex CLI. Compact comprime el contenido de la parte del diálogo de múltiples turnos, pero conserva los mensajes de sistema preinyectados. Este es un recurso clave de la Context Engineering para administrar la ventana de contexto.
\end{importantbox}

\subsection{Carga progresiva de los Skills}

\begin{knowledgebox}{El mecanismo de carga en dos etapas de los Skills}
Los Skills adoptan una estrategia de carga progresiva (Progressive Loading):
\begin{enumerate}
    \item \textbf{Etapa uno (registro)}: en el mensaje de User preinyectado, solo se exponen el nombre y una breve descripción de cada Skill
    \item \textbf{Etapa dos (carga)}: cuando el modelo remoto decide usar un Skill concreto, se ejecuta el comando \texttt{cat} mediante una Function Call, y el contenido completo del archivo skills.md correspondiente a ese Skill se carga en el Context como Function Output
\end{enumerate}
Este diseño evita el desperdicio de Context que resultaría de cargar todos los Skills de una sola vez.
\end{knowledgebox}

\subsection{Resumen de la sección}

El historial de mensajes de Codex crece de manera continua en tiempo de ejecución, y se comprime el contexto mediante la instrucción compact. Los Skills adoptan una carga progresiva: primero se expone un resumen y después se carga el contenido completo según se necesite, lo cual constituye una estrategia refinada de administración del Context.
